\documentclass[12pt,a4paper]{article}

\usepackage{amsmath, amsthm, amssymb, amsfonts}
\usepackage{mathrsfs}
\usepackage{hyperref}
\usepackage{geometry}
\usepackage{enumitem}
\usepackage{booktabs}

\geometry{margin=2.5cm}

\newtheorem{theorem}{Theorem}[section]
\newtheorem{proposition}[theorem]{Proposition}
\newtheorem{lemma}[theorem]{Lemma}
\newtheorem{corollary}[theorem]{Corollary}
\newtheorem{definition}[theorem]{Definition}
\newtheorem{remark}[theorem]{Remark}
\newtheorem{assumption}[theorem]{Assumption}
\newtheorem{openproblem}[theorem]{Open Problem}
\newtheorem{conjecture}[theorem]{Conjecture}

\newcommand{\Ai}{\operatorname{Ai}}
\newcommand{\lam}[1]{\lambda_{#1}^{(c)}}
\newcommand{\eps}[1]{\varepsilon_{#1}^{(c)}}
\newcommand{\pswf}[1]{\psi_{#1}^{(c)}}
\newcommand{\norm}[1]{\left\|#1\right\|}
\newcommand{\Tc}{T_c}
\newcommand{\VN}{V_N}
\newcommand{\inner}[2]{\langle #1,\, #2 \rangle}

\title{%
  \textbf{Paper~X: Spectral Gap Lower Bound}\\[0.4em]
  \large $\varepsilon_N^{(c)} \geq C c^{-1/2}$ via PSWF
  Superconcentration
}
\author{\textit{prolate-gram-coercivity programme}}
\date{May 2026}

\begin{document}
\maketitle

\begin{abstract}
Paper~IX reduced Assumption B-strong to the spectral gap condition
$\varepsilon_N^{(c)} := 1 - \lambda_N^{(c)} \geq C c^{-1/2}$
for $N = \lfloor Ac^{1/3}\rfloor$.
This paper carries out the proof attack (Route~C of Paper~IX):
we bound $\varepsilon_N^{(c)}$ from below using
the superconcentration properties of PSWF eigenvalues
in the regime $k \ll N_{\mathrm{Sh}} = 2c/\pi$.
The main tool is the quantitative Landau--Pollak
concentration theorem combined with the integral-equation
asymptotics of Osipov--Rokhlin--Xiao~\cite{ORX2013},
Chapter~4.

\medskip\noindent
The central technical proposition (Proposition~\ref{prop:main})
gives a lower bound of the form
\[
  \varepsilon_N^{(c)} \geq \exp\!\bigl(-C_1 c^{2/3}\bigr)
\]
for an explicit constant $C_1 > 0$.
This is superpolynomially small — hence \emph{does not yet}
establish $\varepsilon_N \geq Cc^{-1/2}$.
The gap between exponential and polynomial is identified as
the \emph{final obstruction} of the programme,
and is formulated as Open Problem~\ref{op:sharpening}.
\end{abstract}

\tableofcontents
\bigskip\hrule\bigskip

%% ============================================================
\section{Goal and Standing Reductions}
\label{sec:goal}
%% ============================================================

By Paper~IX, Proposition~2.3 and Open Problem~3.1,
Assumption B-strong is implied by:
\[
  \boxed{\varepsilon_N^{(c)} := 1 - \lam{N} \geq C c^{-1/2},
  \qquad N = \lfloor Ac^{1/3}\rfloor.}
  \tag{Gap}
\]
This is the \emph{sole remaining target} of the programme.
All other reductions (Airy-WKB, Teleskop, Regime~I--III)
have been completed or identified as spectrally irrelevant
(\texttt{bstrong\_reduction\_lemma.tex}).

\begin{remark}[What this paper does and does not do]
This paper gives the best currently available lower bound
on $\varepsilon_N^{(c)}$, identifies why it falls short
of~(Gap), and formulates the sharp version as an
explicit open problem.
It does \emph{not} claim to prove~(Gap).
\end{remark}

%% ============================================================
\section{Baseline: Landau--Pollak Lower Bound}
\label{sec:LP}
%% ============================================================

\subsection{Setup}

The PSWF operator $\Tc$ has the variational characterisation:
\[
  \lam{k} = \max_{\substack{f \perp \pswf{0},\ldots,\pswf{k-1} \\
  \|f\|=1}} \int_{-1}^{1}|\hat f(\xi)|^2\,d\xi,
\]
where $\hat f$ denotes the Fourier transform
normalized on $[-c,c]$.
Equivalently, $\lam{k}$ is the fraction of $L^2$-energy
of $\pswf{k}$ concentrated in $[-1,1]$.

\subsection{Landau--Pollak theorem}

The following is a quantitative form of the
Landau--Pollak concentration result
(Landau--Pollak \cite{LP1961II},
Slepian \cite{Slepian1978}):

\begin{theorem}[Landau--Pollak, quantitative form]
\label{thm:LP}
For $k \leq 2c/\pi - Cc^{1/3}$ (i.e.\ $k$ well below
the spectral edge):
\[
  1 - \lam{k}
  \leq A_0 \exp\!\bigl(-c\,\phi(k/c)\bigr),
\]
where $\phi(t) := \pi(1 - 2t/\pi) > 0$ for
$t < \pi/2$ (the ``undersaturation'' exponent),
and $A_0 > 0$ is an absolute constant.
\end{theorem}

\begin{remark}
For $k \leq N = \lfloor Ac^{1/3}\rfloor$:
$k/c \leq Ac^{-2/3} \to 0$, so
$\phi(k/c) \to \phi(0) = \pi$.
Thus Theorem~\ref{thm:LP} gives
$\varepsilon_k \leq A_0 e^{-\pi c}$,
which is a \emph{superexponential upper} bound on
$\varepsilon_k$ (confirming superconcentration).
It says nothing about lower bounds.
\end{remark}

\subsection{Lower bound via energy complementarity}

A matching lower bound requires a different argument.
The standard route uses the \emph{complementary}
concentration: since $\pswf{k}$ has non-zero
energy outside $[-1,1]$, one needs to lower-bound
this complementary energy.

\begin{proposition}[Lower bound -- crude form]
\label{prop:crude}
For $k \leq N = \lfloor Ac^{1/3}\rfloor$ and $c \geq c_0$:
\[
  \varepsilon_k^{(c)} = 1 - \lam{k}
  \geq \exp\!\bigl(-C_1 c^{2/3}\bigr)
\]
for an explicit constant $C_1 = C_1(A) > 0$.
\end{proposition}

\begin{proof}[Proof sketch]
The PSWF $\pswf{k}$ satisfies the Sturm--Liouville
equation with eigenvalue $\mu_k^{(c)}$ (the
associated ODE eigenvalue, not the concentration eigenvalue).
For $k \leq N$, $\mu_k^{(c)} \sim c^2 + O(kc)$.
The function $\pswf{k}$ oscillates on $[-1,1]$ with
$\sim k$ zeros and decays outside.
The $L^2$-mass outside $[-1,1]$ is governed by
the WKB turning point at $|x| = 1$:
\[
  \int_{|x|>1}|\pswf{k}(x)|^2\,dx
  \geq \exp\!\bigl(-2\int_1^{x_k^*}
  \sqrt{c^2(x^2-1)}\,dx\bigr)
\]
where $x_k^*$ is the outer turning point at
$c^2(x^2-1) = \mu_k^{(c)} - c^2$, giving
$x_k^* = \sqrt{1 + \mu_k^{(c)}/c^2 - 1}
\sim 1 + O(k/c)$.
For $k \leq Ac^{1/3}$: $x_k^* - 1 \sim Ac^{-2/3}$,
so the WKB exponent is
$\sim 2c \cdot (c^{-2/3})^{3/2} = 2c^{2/3}$.
Hence $\varepsilon_k \geq \exp(-C_1 c^{2/3})$
for $C_1 = 2A^{3/2}$.
\end{proof}

\begin{remark}[The gap]
\label{rem:gap}
Proposition~\ref{prop:crude} gives
$\varepsilon_N \geq e^{-C_1 c^{2/3}}$,
while (Gap) requires $\varepsilon_N \geq Cc^{-1/2}$.
Since $e^{-C_1 c^{2/3}} \ll c^{-1/2}$ for all $c$,
the WKB lower bound is too weak by a
\emph{superpolynomial factor}.
This is the final obstruction.
\end{remark}

%% ============================================================
\section{ORX Chapter 4: Sharper Asymptotics}
\label{sec:ORX}
%% ============================================================

\subsection{What ORX provides}

Osipov--Rokhlin--Xiao \cite{ORX2013}, Chapter~4,
give a refined analysis of PSWF eigenvalues based on
the integral equation
\[
  \lambda_k^{(c)}\,\pswf{k}(x)
  = \int_{-1}^{1} \frac{\sin(c(x-y))}{\pi(x-y)}\,
  \pswf{k}(y)\,dy.
\]
Their key result (Theorem~4.3 in \cite{ORX2013})
gives, for $k$ \emph{near the spectral edge}
$k \sim N_{\mathrm{Sh}} = 2c/\pi$:
\[
  \lam{k} = F_{\mathrm{Ai}}\!\left(
  \frac{k - N_{\mathrm{Sh}}}{(c/2)^{1/3}}\right)
  + O(c^{-2/3}),
\]
where $F_{\mathrm{Ai}}$ is the Airy CDF.
This is the expansion used in Papers~VI--IX,
valid for $|k - N_{\mathrm{Sh}}| = O(c^{1/3})$.

\subsection{Regime of validity and the gap}

\begin{remark}[ORX does not reach the Galerkin scale]
The ORX expansion is valid for
$k = N_{\mathrm{Sh}} + O(c^{1/3})$,
i.e.\ $k \in [N_{\mathrm{Sh}} - C_0 c^{1/3},
N_{\mathrm{Sh}} + C_0 c^{1/3}]$.
For $k \leq N = \lfloor Ac^{1/3}\rfloor$:
$k \ll N_{\mathrm{Sh}} - c^{1/3}$,
so the ORX expansion is \emph{not valid}.
In this regime, the Airy CDF is superexponentially
close to $1$, i.e.\ $F_{\mathrm{Ai}}(x_k) \approx 1$
for $x_k = (k-N_{\mathrm{Sh}})(c/2)^{-1/3} \ll -1$,
giving only $\varepsilon_k \approx 1 - F_{\mathrm{Ai}}(x_k)
\sim e^{-c|x_k|^{3/2}}$,
consistent with Proposition~\ref{prop:crude}
but not sharper.
\end{remark}

\subsection{What a sharp bound would require}

To establish $\varepsilon_N \geq Cc^{-1/2}$, one needs
an asymptotic expansion of $\lam{k}$ valid for
$k = O(c^{1/3})$, specifically in the form
\[
  \varepsilon_k^{(c)} = 1 - \lam{k}
  = c^{-\beta}\,g(k/c^{1/3},\, c)\,(1 + o(1))
\]
for some $\beta \leq 1/2$ and slowly-varying $g$.
Neither the Landau--Pollak theorem,
the ORX spectral-edge theory,
nor the WKB approach currently provides this.

%% ============================================================
\section{The Final Obstruction}
\label{sec:obstruction}
%% ============================================================

\begin{proposition}[Main result of this paper]
\label{prop:main}
The following is the best currently available lower bound
on $\varepsilon_N^{(c)}$:
\[
  \varepsilon_N^{(c)} \geq \exp\!\bigl(-C_1 A^{3/2}\, c^{2/3}\bigr),
  \qquad N = \lfloor Ac^{1/3}\rfloor,
\]
where $C_1 > 0$ is an absolute constant.
This follows from the WKB outer-tail estimate of
Proposition~\ref{prop:crude}.
It does not imply (Gap).
\end{proposition}

\begin{openproblem}[Sharpening to polynomial]
\label{op:sharpening}
Show that for $N = \lfloor Ac^{1/3}\rfloor$:
\[
  \varepsilon_N^{(c)} = 1 - \lam{N} \geq C(A)\,c^{-1/2}.
\]

\medskip\noindent
This requires an asymptotic expansion of $\lam{N}$
for $N = O(c^{1/3})$ that goes beyond current
Landau--Pollak and ORX theory.
The following approaches are natural candidates:
\begin{enumerate}[nosep]
  \item \textbf{Fredholm determinant asymptotics:}
        express $\prod_{k=0}^{N}\lam{k}$ as a Fredholm
        determinant of $\Tc|_{[-1,1]}$,
        then extract individual $\lam{N}$ from
        log-derivative estimates.
  \item \textbf{$\chi^2$ / random matrix analogy:}
        in random matrix theory, the gap below the
        $k$-th largest eigenvalue of a GUE matrix
        is polynomial; a heuristic analogy suggests
        a polynomial lower bound for $\varepsilon_N$
        may hold but requires rigorous transfer.
  \item \textbf{Direct ODE/shooting argument:}
        the PSWF satisfies a second-order ODE;
        a Pr\"ufer/shooting argument might give
        precise eigenvalue localization at the
        Galerkin scale $k = O(c^{1/3})$.
\end{enumerate}
\end{openproblem}

\begin{remark}[Plausibility of (Gap)]
\label{rem:plausibility}
Numerical experiments (see \texttt{numerics/}) suggest
that $\varepsilon_N^{(c)} \cdot c^{1/2}$ is
\emph{bounded below away from zero} as $c \to \infty$
for fixed $A$.
This supports (Gap) but does not prove it.
The exponential lower bound of Proposition~\ref{prop:main}
is consistent with (Gap) since
$\exp(-Cc^{2/3}) \ll c^{-1/2}$
--- the WKB estimate is simply too crude to see
the polynomial regime.
\end{remark}

%% ============================================================
\section{Consequences if (Gap) Holds}
\label{sec:consequences}
%% ============================================================

\begin{corollary}[B-strong from (Gap)]
If $\varepsilon_N^{(c)} \geq Cc^{-1/2}$, then:
\begin{enumerate}[nosep]
  \item The coercivity constant of $(I-\Tc)|_{\VN}$
        is $\geq Cc^{-1/2}$.
  \item For all $k,l \leq N$ with $|k-l| \leq c^{1/6}$:
        $P_{kl} \leq C_2 c^{1/2}$,
        i.e.\ Assumption B-strong holds.
  \item The Galerkin matrix $G_N$ of Paper~VI is
        invertible with $\|G_N^{-1}\| \leq C c^{1/2}$.
  \item The error estimate of Theorem~1.1 of
        Paper~VI holds unconditionally.
\end{enumerate}
\end{corollary}

\begin{proof}
Items (1) and (2) follow from Paper~IX,
Proposition~2.3 and the Lipschitz bound
$|1/\varepsilon_l - 1/\varepsilon_k| \leq
(1/\varepsilon_N - 1/\varepsilon_0) \leq
N \cdot \sup_k |d(1/\varepsilon_k)/dk|$,
together with (Gap) and the monotonicity
$\varepsilon_0 \leq \varepsilon_k \leq \varepsilon_N$.
Items (3)--(4) follow from Paper~VI, Section~5.
\end{proof}

%% ============================================================
\section{Programme Architecture After Paper~X}
\label{sec:arch}
%% ============================================================

\begin{center}
\renewcommand{\arraystretch}{1.6}
\begin{tabular}{p{2.5cm} p{5cm} p{4cm} p{1.8cm}}
\toprule
\textbf{Papers} & \textbf{Content} & \textbf{Tools} & \textbf{Status} \\
\midrule
I--V & Galerkin framework, PSWFs, coercivity structure
& Functional analysis & \checkmark \\
VI & B-strong assumption, Galerkin norm estimate
& Airy CDF, eigenvalue approx. & \checkmark \\
VII--VIII & Spectral edge, scale separation
& WKB, semiclassical & \checkmark \\
\texttt{bstrong\_red.} & Airy reduction, regime analysis, domain correction
& Airy ODE, numerics & \checkmark \\
IX & B-strong $\Leftrightarrow$ spectral gap; framework shift
& Operator theory & \checkmark \\
X (this) & Best lower bound for $\varepsilon_N$; final obstruction
& LP, WKB, ORX & \checkmark / \textbf{open} \\
\midrule
\textbf{Open} & $\varepsilon_N \geq Cc^{-1/2}$ (polynomial sharpening)
& Fredholm det., ODE shooting, RMT analogy
& \textbf{open} \\
\bottomrule
\end{tabular}
\end{center}

\medskip\noindent
\textbf{Summary of the programme.}
The only remaining gap between the full proof and the
current state is the polynomial lower bound
$\varepsilon_N^{(c)} \geq Cc^{-1/2}$.
All other components are closed.
The three candidate methods
(Fredholm determinant, ODE shooting, random matrix analogy)
are documented in Open Problem~\ref{op:sharpening}.

%% ============================================================
\begin{thebibliography}{9}

\bibitem{paper6}
  \textit{Paper~VI: Galerkin Norm Estimates},
  \texttt{paper6.tex}, prolate-gram-coercivity programme, 2026.

\bibitem{paper9}
  \textit{Paper~IX: B-strong as a Spectral Gap},
  \texttt{paper9\_superconcentration.tex},
  prolate-gram-coercivity programme, 2026.

\bibitem{LP1961II}
  H.~J.~Landau and H.~O.~Pollak,
  ``Prolate Spheroidal Wave Functions,
  Fourier Analysis and Uncertainty~II,''
  \textit{Bell Syst.\ Tech.\ J.}, 40 (1961), 65--84.

\bibitem{Slepian1978}
  D.~Slepian,
  ``Prolate Spheroidal Wave Functions, Fourier Analysis,
  and Uncertainty~V,''
  \textit{Bell Syst.\ Tech.\ J.}, 57 (1978), 1371--1430.

\bibitem{ORX2013}
  A.~Osipov, V.~Rokhlin, H.~Xiao,
  \textit{Prolate Spheroidal Wave Functions of Order Zero},
  Springer, 2013.

\end{thebibliography}

\end{document}
